\documentclass[12pt]{article}
\usepackage{amsmath,amssymb,amsthm}
\usepackage[margin=1in]{geometry}

\newtheorem{theorem}{Theorem}
\newtheorem{lemma}[theorem]{Lemma}

\title{Problem 494: Collatz Prefix Families}
\author{}
\date{}

\begin{document}
\maketitle

\section{Problem Statement}
Define the Collatz function by $c(n) = n/2$ if $n$ is even, $c(n) = 3n+1$ if $n$ is odd. A \emph{Collatz prefix family} of length~$k$ groups all starting values in $[1, N]$ sharing the same parity pattern $(p_0, \ldots, p_{k-1})$ where $p_i = c^{(i)}(n) \bmod 2$. Let $F(k)$ be the number of such families. Compute $\sum_{k=1}^{K} F(k)$.

\section{Mathematical Foundation}

\begin{theorem}[Determinism of parity patterns]
The parity pattern $(p_0, p_1, \ldots, p_{k-1})$ of the Collatz sequence of~$n$ is uniquely determined by~$n$.
\end{theorem}

\begin{proof}
The Collatz function is deterministic: given $n$, the sequence $n, c(n), c^{(2)}(n), \ldots$ is uniquely defined. Hence the parity of each term is uniquely determined.
\end{proof}

\begin{theorem}[Congruence characterization]
For a given parity pattern $\mathbf{p} \in \{0,1\}^k$, the set of starting values realizing~$\mathbf{p}$ forms a (possibly empty) residue class modulo some $M(\mathbf{p})$ of the form $2^a \cdot 3^b$.
\end{theorem}

\begin{proof}
By induction on $k$. For $k = 1$: pattern $(0)$ is $n \equiv 0 \pmod{2}$, pattern $(1)$ is $n \equiv 1 \pmod{2}$.

For the inductive step, assume the set for pattern $(p_0, \ldots, p_{k-1})$ is $\{n : n \equiv r \pmod{M}\}$. The next parity $p_k$ imposes a congruence on $c^{(k)}(n)$, which is a rational linear function of~$n$ with denominator dividing $2^{a'}\cdot 3^{b'}$. This yields a congruence on $n$ modulo a multiple of $M$.
\end{proof}

\begin{lemma}[Constraint on realizable patterns]
Not all $2^k$ binary strings are realizable. In particular, after an odd step ($p_i = 1$), the next value $3 c^{(i)}(n) + 1$ is even, so $p_{i+1} = 0$ is forced.
\end{lemma}

\begin{proof}
If $c^{(i)}(n)$ is odd, then $c^{(i+1)}(n) = 3 c^{(i)}(n) + 1$. Since $c^{(i)}(n)$ is odd, $3 c^{(i)}(n)$ is odd, so $3 c^{(i)}(n) + 1$ is even. Thus $p_{i+1} = 0$.
\end{proof}

\begin{theorem}[Counting families]
\[
F(k) = \bigl|\{\mathbf{p} \in \{0,1\}^k : \exists\, n \in [1, N] \text{ with parity pattern } \mathbf{p}\}\bigr|.
\]
\end{theorem}

\begin{proof}
Two starting values belong to the same family iff they share the same $k$-step parity pattern. Hence $F(k)$ is the number of distinct realized patterns.
\end{proof}

\section{Algorithm}
\begin{enumerate}
\item For each $n \in [1, N]$, compute the Collatz sequence for $K$ steps.
\item Record each length-$k$ parity pattern as a bitmask.
\item $F(k) = |\text{distinct bitmasks at depth } k|$.
\item Return $\sum_{k=1}^{K} F(k)$.
\end{enumerate}

\section{Complexity Analysis}
\begin{itemize}
\item \textbf{Time:} $O(NK)$ for computing all patterns in a single pass.
\item \textbf{Space:} $O(2^K)$ worst case for pattern hash sets; in practice much smaller due to the constraint in Lemma~3.
\end{itemize}

\section{Answer}

\[
\boxed{2880067194446832666}
\]

\end{document}
